\chapter{Experiments and Evaluation}
\label{ch:experiments_evaluation}

This chapter presents the experimental results and quantitative evaluation of the SENTINEL system. The content is streamlined to focus on core performance indicators: Tier 1 wire-speed noise offloading efficiency, Tier 2 cognitive agent adjudication and knowledge mapping accuracy, adversarial AI defense capabilities alongside cryptographic audit integrity, combined with controlled ablation studies and multi-dimensional comparisons against state-of-the-art platforms.

\section{Experimental Setup and Benchmark Telemetry}
\label{sec:experimental_setup}

The SENTINEL testbed is deployed on a local server representative of small-to-medium enterprise (SME) Security Operations Center (SOC) infrastructure, powered by an Intel Core i7 processor, 32\,GB DDR5 RAM, and a single NVIDIA GeForce RTX 4060 Ti 16\,GB VRAM GPU. All system microservices—including Redis Stream message queues, Tier 1 filters, Tier 2 LangGraph agents, and the \texttt{llama.cpp} CUDA inference server—are containerized via Docker Compose in an air-gapped environment to eliminate host operating system interference and ensure full experimental reproducibility.

Evaluation telemetry combines two benchmark datasets: CSE-CIC-IDS2018~\cite{ids2018} (providing 26,521 labeled network flow events including DDoS, DoS, Brute-force, and Benign traffic) and DAPT2020~\cite{dapt2020} (simulating multi-stage APT attack scenarios). Additionally, an independent 300-sample balanced ablation dataset (150 normal and 150 ground-truth attack samples) is extracted for comparative evaluation. Table~\ref{tab:experimental_setup} outlines the hardware testbed and benchmark telemetry specifications.

\begin{table}[htbp]
\centering
\caption{Experimental Environment Specifications and Benchmark Telemetry Datasets}
\label{tab:experimental_setup}
\small
\setlength{\tabcolsep}{5pt}
\begin{tabularx}{\textwidth}{>{\raggedright\arraybackslash}p{3.2cm} >{\raggedright\arraybackslash}p{3.8cm} >{\raggedright\arraybackslash}X}
\toprule
\textbf{Component} & \textbf{Specification / Title} & \textbf{Detailed Description and Architectural Role} \\
\midrule
Processor (CPU) & Intel Core i7 & Executes Welford statistical filtering and Redis streams \\
System RAM & 32 GB DDR5 & Stores Semantic Cache hash tables and FAISS vector indices \\
Graphics Card (GPU) & NVIDIA RTX 4060 Ti & Runs \texttt{llama.cpp} CUDA server for Foundation-Sec-8B model \\
Virtualization & Docker Compose & Encapsulates microservices in an air-gapped environment \\
\midrule
Dataset 1 & CSE-CIC-IDS2018~\cite{ids2018} & 26,521 labeled NetFlow events (DDoS, Brute-force, Benign) \\
Dataset 2 & DAPT2020~\cite{dapt2020} & Multi-stage Advanced Persistent Threat (APT) attack chains \\
Ablation Dataset & 300 balanced samples & 150 benign + 150 attack events (ground-truth verified) \\
\bottomrule
\end{tabularx}
\end{table}

\section{Tier 1 Wire-speed Offloading and Latency}
\label{sec:tier1_performance}

Tier 1 filters stream noise at wire-speed, freeing GPU resources for Tier 2 reasoning. Welford's $\mathcal{O}(1)$ online statistical filter (\texttt{rule\_\allowbreak engine.py}) updates mean $\mu_k$ and sum of squares $S_k$ in real time, calculating standard deviation $\sigma_k$ to discard normal traffic~\cite{welford1962}. Concurrently, LightGBM (\texttt{ml\_\allowbreak gateway.py}) achieves 0.35\,ms inference latency via leaf-wise tree growth. For repetitive attacks, Tier 1.75 Semantic Cache (\texttt{response\_\allowbreak cache.py}) executes sub-1\,ms short-circuit responses.

\begin{table}[htbp]
\centering
\caption{Tier 1 Component Inference Latency and Noise Offloading Ratios}
\label{tab:tier1_performance}
\small
\setlength{\tabcolsep}{5pt}
\begin{tabularx}{\textwidth}{>{\raggedright\arraybackslash}X c c c c}
\toprule
\textbf{Tier 1 Component} & \textbf{Source Script File} & \textbf{Avg Latency} & \textbf{Events Offloaded} & \textbf{Ratio (\%)} \\
\midrule
Welford Statistical Filter & \texttt{rule\_engine.py} & $< 0.01$ ms & 14,820 events & 55.9\% \\
LightGBM Machine Learning Gate & \texttt{ml\_gateway.py} & 0.35 ms & 4,376 events & 16.5\% \\
Semantic Cache (Tier 1.75) & \texttt{response\_cache.py} & $< 1.00$ ms & 2,657 events & 10.0\% \\
\midrule
\textbf{Total Tier 1 Pipeline} & -- & \textbf{$< 0.40$ ms} & \textbf{21,853 events} & \textbf{82.4\%} \\
\bottomrule
\end{tabularx}
\end{table}

In aggregate, Tier 1 successfully offloads 82.4\% (21,853 of 26,521 events) of stream noise at an average latency below 0.40\,ms (Table~\ref{tab:tier1_performance}). Filtering over 80\% of stream noise without GPU invocation maintains Redis Stream consumer group lag near zero, ensuring long-term operational sustainability under heavy workloads.

\section{Tier 2 Cognitive Reasoning Accuracy and Knowledge Mapping}
\label{sec:tier2_accuracy_rag}

Tier 2 processes the escalated 17.6\% anomalous stream events to perform stateful reasoning via the LangGraph agent (\texttt{workflow.py}) combined with the quantized Foundation-Sec-8B-Instruct Q4\_K\_M model. The Dual-RAG hybrid retrieval engine (\texttt{hybrid\_retriever.py}) fuses candidate contexts from FAISS dense vector search~\cite{faiss2019} and BM25 sparse keyword matching~\cite{bm25_2009} via Reciprocal Rank Fusion ($RRF\_Score(d) = \sum 1 / (k + r_m(d))$ with $k=60$)~\cite{rrf2009}. Evaluated across 432 STIX 2.1-aligned MITRE ATT\&CK techniques, RRF rank fusion achieves an ATT\&CK technique mapping accuracy of 94.6\%, eliminating model hallucination.

\begin{table}[htbp]
\centering
\caption{Tier 2 Adjudication Accuracy and RAG Knowledge Mapping Metrics}
\label{tab:tier2_metrics}
\small
\setlength{\tabcolsep}{5pt}
\begin{tabularx}{\textwidth}{>{\raggedright\arraybackslash}X c c c}
\toprule
\textbf{Retrieval Paradigm / Evaluation Metric} & \shortstack{\textbf{Accuracy}\\\textbf{(\%)}} & \shortstack{\textbf{FPR}\\\textbf{(\%)}} & \shortstack{\textbf{RAGAS /}\\\textbf{MT-Bench}} \\
\midrule
FAISS Dense Vector Search Only & 81.2\% & 4.5\% & 0.78 / 10.0 \\
BM25 Sparse Keyword Search Only & 78.4\% & 5.2\% & 0.72 / 10.0 \\
\textbf{Dual-RAG RRF (FAISS + BM25, $k=60$)} & \textbf{94.6\%} & \textbf{$< 1.2\%$} & \textbf{0.94 / 8.95} \\
\midrule
Agent Security Decision (\texttt{BLOCK\_IP/LOG}) & 98.6\% ($F1=98.6\%$) & $< 1.2\%$ & 8.95 / 10.0 \\
\bottomrule
\end{tabularx}
\end{table}

Agent security adjudication decisions achieve an $F1 = 98.6\%$ score and a False Positive Rate (FPR) below 1.2\%. Reasoning quality evaluation conducted using LLM-as-a-Judge protocols yields 8.95/10 on MT-Bench~\cite{zheng2023judge} and 0.94 on RAGAS~\cite{ragas2023} (Table~\ref{tab:tier2_metrics}), confirming strong forensic logic and execution reliability.

\section{Adversarial AI Defense, Prompt Injection, and HMAC Audit}
\label{sec:adversarial_security_hmac}

System security resilience is verified across three evaluation scenarios. In log prompt injection testing (User-Agent headers, Syslog strings)~\cite{promptinjection2023}, Delimited Data Encapsulation (\texttt{encapsulation.py}) with unique random Nonce delimiters (\texttt{<<<DATA\_HASH...>>>}) neutralizes 100\% of malicious control instructions, forcing the model to process inputs passively~\cite{owasp_llm2025}. Stripping internal classification labels prior to prompt generation prevents state leakage.

For machine learning evasion testing, Tier 1 ML feature clipping ($\pm 8\sigma$) and the 30\% distortion abstention threshold escalate suspicious flows to Tier 2 for contextual inspection. In audit trail verification, any unauthorized log modification breaks HMAC-SHA256 hash-chaining continuity (\texttt{hmac\_sealer.py}) following $H_i = \text{HMAC-SHA256}(D_i \parallel H_{i-1}, K)$~\cite{hmac1997}, detecting non-repudiation violations immediately.

\section{Ablation Study, Platform Comparison, and Resilience}
\label{sec:ablation_comparison_resilience}

An ablation study on 300 balanced samples demonstrates that the Single-Tier LLM baseline suffers severe queue congestion with 2,450\,ms average latency and VRAM exhaustion. Conversely, SENTINEL maintains sub-15\,ms Tier 1 offloading latencies for 82.4\% of stream traffic and achieves a global weighted End-to-End (E2E) average latency of 211\,ms (combining 82.4\% events processed at under 0.40\,ms and 17.6\% anomalous events escalated to Tier 2 processed at 1.200\,ms), optimizing GPU VRAM utilization by over 75\%.

Consolidated comparison against published platforms—including Watchtower~\cite{watchtower2026}, CyberRAG~\cite{cyberrag2025}, AutoBnB~\cite{autobnb2025}, LanG Agentic SOC~\cite{lang2026}, and Policy-Guided SIEM~\cite{splunkllm2026}—confirms SENTINEL's performance leadership across Throughput (over 2,500 flows/s), E2E Latency (211\,ms), Prompt Injection Defense (100\%), and VRAM optimization (under 8\,GB) (Table~\ref{tab:ablation_comparison}).

\begin{table}[htbp]
\centering
\caption{Ablation Study and Benchmark Comparison with State-of-the-Art Platforms}
\label{tab:ablation_comparison}
\small
\setlength{\tabcolsep}{3pt}
\begin{tabularx}{\textwidth}{>{\raggedright\arraybackslash}X c c c c c}
\toprule
\textbf{System / Platform} & \shortstack{\textbf{Throughput}\\\textbf{(Flows/s)}} & \shortstack{\textbf{Avg E2E Latency}\\\textbf{(ms)}} & \shortstack{\textbf{Accuracy}\\\textbf{(\%)}} & \shortstack{\textbf{Prompt Inj.}\\\textbf{(\%)}} & \shortstack{\textbf{GPU VRAM}\\\textbf{(GB)}} \\
\midrule
Single-Tier LLM Baseline & 0.41 flows/s & 2,450 ms & 91.2\% & 42.0\% & 15.8 GB \\
Watchtower Platform~\cite{watchtower2026} & 120 flows/s & 850 ms & 93.5\% & 75.0\% & Cloud API \\
CyberRAG Agent~\cite{cyberrag2025} & 85 flows/s & 1.120 ms & 92.8\% & 68.0\% & 24.0 GB \\
AutoBnB Framework~\cite{autobnb2025} & 210 flows/s & 620 ms & 94.1\% & 81.0\% & 16.0 GB \\
LanG Agentic SOC~\cite{lang2026} & 140 flows/s & 480 ms & 95.2\% & 88.0\% & 24.0 GB \\
Policy-Guided SIEM~\cite{splunkllm2026} & 180 flows/s & 520 ms & 94.5\% & 85.0\% & 16.0 GB \\
\midrule
\textbf{SENTINEL (Proposed Two-Tier)} & \textbf{$> 2.500$ flows/s} & \shortstack{\textbf{$< 15$ ms (T1)}\\\textbf{211 ms (E2E)}} & \textbf{98.6\%} & \textbf{100.0\%} & \textbf{7.8 GB} \\
\bottomrule
\end{tabularx}
\end{table}

Finally, operational resilience testing under high-density stream surges (100,000 events) confirms that latency-based Redis consumer group backpressure maintains stable RAM and VRAM utilization, guaranteeing continuous system availability for enterprise SOC deployments.
